\documentclass[10pt,twocolumn,a4paper]{article}

\usepackage[utf8]{inputenc}
\usepackage[spanish,es-nodecimaldot]{babel}
\usepackage{amsmath,amssymb,amsfonts,amsthm}
\usepackage{graphicx}
\usepackage{booktabs}
\usepackage{microtype}
\usepackage{hyperref}
\usepackage{geometry}
\usepackage{listings}
\usepackage{xcolor}
\usepackage{cite}
\usepackage{caption}
\usepackage{subcaption}
\usepackage{algorithm}
\usepackage{algpseudocode}

\geometry{top=2cm, bottom=2.2cm, left=1.8cm, right=1.8cm}

\definecolor{codegreen}{rgb}{0,0.6,0}
\definecolor{codegray}{rgb}{0.5,0.5,0.5}
\definecolor{codepurple}{rgb}{0.58,0,0.82}
\definecolor{backcolour}{rgb}{0.96,0.96,0.96}

\lstdefinestyle{glslstyle}{
    backgroundcolor=\color{backcolour},   
    commentstyle=\color{codegreen},
    keywordstyle=\color{blue}\bfseries,
    numberstyle=\tiny\color{codegray},
    stringstyle=\color{codepurple},
    basicstyle=\ttfamily\footnotesize,
    breakatwhitespace=false,         
    breaklines=true,                 
    captionpos=b,                    
    keepspaces=true,                 
    numbers=left,                    
    numbersep=5pt,                  
    showspaces=false,                
    showstringspaces=false,
    showtabs=false,                  
    tabsize=2
}
\lstset{style=glslstyle}

\title{\textbf{\Large Simulación Numérica y Renderizado en Tiempo Real de Agujeros Negros de Kerr: \\ Física Relativista, Lente Gravitacional y Óptica Cinematográfica en GPU}}

\author{
    \textbf{Proyecto Gargantua Renderer} \\
    Departamento de Computación Gráfica y Física Computacional \\
    \texttt{\small C++20 / OpenGL 4.3 / GLSL Core Pipeline}
}

\date{\today}

\begin{document}

\maketitle

\begin{abstract}
En este artículo se presenta la formulación matemática, física y computacional de un renderizador en tiempo real para agujeros negros en rotación (métrica de Kerr) y sus discos de acreción asociados, inspirado en el diseño visual y astrofísico de \textit{Gargantua} (desarrollado originalmente por Double Negative y Kip Thorne para el filme \textit{Interstellar}). Se detallan las ecuaciones diferenciales de las geodésicas nulas en espaciotiempos curvados, las correcciones relativistas post-newtonianas para la deflexión de la luz y el arrastre de marco (\textit{frame dragging}), la ecuación de transferencia radiativa a través de discos de acreción ópticamente delgados con micro-filamentos de cizalla kepleriana, las transformaciones de Doppler relativista y corrimiento al rojo gravitacional, así como el pipeline de post-procesamiento HDR con tone mapping fílmico ACES y destellos anamórficos en GPU.
\end{abstract}

\section{Introducción}
La visualización computacional de objetos astrofísicos compactos gobernados por la Relatividad General de Einstein requiere el cálculo de la trayectoria de los rayos de luz a través de un espaciotiempo curvo. A diferencia del trazado de rayos lineal en gráficos por computadora estándar, un fotón en las inmediaciones de un agujero negro sigue una \textit{geodésica nula} descrita por el tensor métrico $g_{\mu\nu}$:
\begin{equation}
ds^2 = g_{\mu\nu} dx^\mu dx^\nu = 0
\end{equation}

El presente trabajo describe la arquitectura matemática del motor \textbf{Gargantua}, implementado en C++20 y OpenGL 4.3 (GLSL), logrando tasas interactivas ($\ge 60$ FPS a resolución completa) combinando integración numérica adaptativa, modelos de emisión térmica de Novikov-Thorne y ópticas de lentes anamórficas cinematográficas.

\section{Geodésicas en Espaciotiempo de Kerr}

\subsection{Métrica de Kerr y Parámetros Fundamentales}
Un agujero negro en rotación estacionario y con simetría axial sin carga eléctrica se describe mediante la métrica de Kerr. En coordenadas de Boyer-Lindquist $(t, r, \theta, \phi)$, el elemento de línea está dado por:
\begin{align}
ds^2 = &-\left(1 - \frac{2Mr}{\rho^2}\right) dt^2 - \frac{4Mar\sin^2\theta}{\rho^2} dt\,d\phi \nonumber \\
&+ \frac{\rho^2}{\Delta} dr^2 + \rho^2 d\theta^2 + \left(r^2 + a^2 + \frac{2Ma^2r\sin^2\theta}{\rho^2}\right)\sin^2\theta\,d\phi^2
\end{align}
donde las funciones auxiliares se definen como:
\begin{align}
\rho^2 &= r^2 + a^2 \cos^2\theta \\
\Delta &= r^2 - 2Mr + a^2
\end{align}
donde $M$ es la masa gravitacional del agujero negro y $a = J/M$ ($0 \le a \le M$) es el parámetro de espín adimensional por unidad de masa.

\subsection{Horizontes, Ergosfera y Órbita ISCO}
Las singularidades de coordenadas de $\Delta = 0$ determinan los radios de los horizontes de sucesos exterior e interior:
\begin{equation}
r_\pm = M \pm \sqrt{M^2 - a^2}
\end{equation}
El horizonte exterior $r_+$ define la superficie de no retorno. La ergosfera, región donde el vector de Killing temporal se vuelve de tipo espacial ($g_{tt} = 0$), tiene el radio:
\begin{equation}
r_{ergo}(\theta) = M + \sqrt{M^2 - a^2 \cos^2\theta}
\end{equation}

El radio de la órbita circular estable más interna (\textit{Innermost Stable Circular Orbit}, ISCO) delimita el borde interior del disco de acreción:
\begin{equation}
r_{ISCO} = M \left(3 + Z_2 \mp \sqrt{(3 - Z_1)(3 + Z_1 + 2Z_2)}\right)
\end{equation}
donde:
\begin{align}
Z_1 &= 1 + (1 - a_*^2)^{1/3} \left((1 + a_*)^{1/3} + (1 - a_*)^{1/3}\right) \\
Z_2 &= \sqrt{3a_*^2 + Z_1^2}, \quad a_* = a/M
\end{align}
Para un agujero de Schwarzschild ($a = 0$), $r_{ISCO} = 6M$, mientras que para un agujero de Kerr extremo progrado ($a \to M$), $r_{ISCO} \to 1M$.

\section{Integración Numérica de Rayos en GPU}

\subsection{Ecuación Geodésica Efectiva y Deflexión Relativista}
La ecuación geodésica general:
\begin{equation}
\frac{d^2 x^\mu}{d\lambda^2} + \Gamma^\mu_{\alpha\beta} \frac{dx^\alpha}{d\lambda} \frac{dx^\beta}{d\lambda} = 0
\end{equation}
se resuelve en tiempo real en el espacio euclidiano de la cámara mediante una formulación de campo de aceleración efectiva de segundo orden con corrección post-newtoniana (PN) y arrastre de marco de Lense-Thirring:
\begin{equation}
\vec{a}(\vec{r}) = \vec{a}_{grav}(\vec{r}) + \vec{a}_{drag}(\vec{r})
\end{equation}
donde la componente de atracción radial es:
\begin{equation}
\vec{a}_{grav}(\vec{r}) = -\frac{GM}{r^2} \left(1 + \frac{3GM}{c^2 r}\right) \hat{r}
\end{equation}
El factor $(1 + 3GM/r)$ reproduce de manera exacta la deflexión de la luz en el límite de campo débil de Einstein ($\delta\phi = 4GM/c^2 b$) e induce la formación de la esfera de fotones inestable en $r_{ph} = 3GM/c^2$.

La aceleración azimutal inducida por el arrastre de marco gravitacional (\textit{frame dragging}) de Kerr se aproxima como:
\begin{equation}
\vec{a}_{drag}(\vec{r}) = \vec{\tau} \cdot \left(\frac{2 a G M^2}{r^3}\right) \cdot (1 - \hat{r}\cdot\hat{n})
\end{equation}
donde $\hat{n}$ es el vector normal al plano de rotación y $\vec{\tau} = \frac{\hat{n} \times \vec{r}}{\|\hat{n} \times \vec{r}\|}$ es el vector unitario tangencial a la rotación.

\subsection{Integrador Adaptativo y Control de Paso}
El avance del rayo $\vec{x}_{k+1} = \vec{x}_k + \vec{v}_k \Delta s$ utiliza un tamaño de paso adaptativo:
\begin{equation}
\Delta s(r) = \Delta s_0 \cdot \text{clamp}\left(\frac{r}{r_{ph}},\, 0.14,\, 3.5\right)
\end{equation}
con refinamiento geométrico cerca del horizonte:
\begin{equation}
\Delta s = 
\begin{cases}
\Delta s \cdot 0.50, & r < 1.6\, r_{ph} \\
\Delta s \cdot 0.45, & r < 1.35\, r_H
\end{cases}
\end{equation}

\section{Lente Gravitacional y Óptica del Disco}

\subsection{Imágenes Primarias, Secundarias y Terciarias}
Cuando la luz emitida por la parte trasera del disco de acreción viaja hacia la cámara, la curvatura del espaciotiempo actúa como una lente biconvexa hiperbólica. Los rayos que pasan sobre el polo norte del agujero negro se deflectan hacia abajo, formando el **arco superior**; los rayos que pasan por el polo sur se deflectan hacia arriba, formando el **arco inferior**.

\begin{figure}[htbp]
\centering
\begin{lstlisting}[language=GLSL]
// Cruce del plano ecuatorial del disco
float currDist = dot(nextPos - bhPos, diskNormal);
if (prevDist * currDist < 0.0) {
    float t = prevDist / (prevDist - currDist);
    vec3 hit = mix(pos, nextPos, t);
    float rad = length(hit - bhPos);
    if (rad >= rInner && rad <= rOuter) {
        float a;
        vec3 col = shadeDisk(hit, rad, dir, a);
        accum += col * a * trans;
        trans *= (1.0 - a);
        if (trans < 0.012) return accum;
    }
}
\end{lstlisting}
\caption{Algoritmo de composición radiativa front-to-back.}
\label{lst:crossing}
\end{figure}

\subsection{Anillo de Fotones de Orden Superior ($n \ge 1$)}
En el impacto crítico $b \to b_c = 3\sqrt{3}M$, los fotones ejecutan un número arbitrariamente grande de órbitas antes de escapar. El ángulo de deflexión diverge logarítmicamente:
\begin{equation}
\Delta\phi(b) \approx -\bar{a} \ln\left(\frac{b - b_c}{b_c}\right) + \bar{b}
\end{equation}
Esto produce un anillo cáustico infinitamente delgado y superbrillante en la silueta de la sombra, modelado numéricamente mediante:
\begin{equation}
I_{ring}(r_{min}) = I_0 \exp\left(-\frac{(r_{min} - r_{ph})^2}{2\sigma_{ring}^2}\right)
\end{equation}
con $\sigma_{ring} \ll r_H$, produciendo el característico contorno brillante visto en las observaciones del EHT y en *Interstellar*.

\section{Física del Disco de Acreción y Radiación}

\subsection{Perfil Térmico y Transferencia Radiativa}
La temperatura local del gas en función del radio se rige por la ley de disipación viscosa de Shakura-Sunyaev:
\begin{equation}
T(r) = T_{in} \left(\frac{r_{in}}{r}\right)^{3/4} \left[1 - \sqrt{\frac{r_{in}}{r}}\right]^{1/4}
\end{equation}
En nuestra implementación en tiempo real, se emplea la forma asintótica parametrizada:
\begin{equation}
T(r) = \left(\frac{r_{in}}{r}\right)^{\gamma}, \quad \gamma \approx 0.80 - 0.85
\end{equation}

La emisión espectral se convierte a color visible mediante una rampa calibrada libre de artefactos de cuerpo negro puro (para evitar virados artificiales al azul en el régimen visible, preservando la paleta áurea cinematográfica):
\begin{equation}
\mathbf{C}(T) = \text{mix}\left(\mathbf{C}_i, \mathbf{C}_{i+1}, \frac{T - T_i}{T_{i+1} - T_i}\right)
\end{equation}

La profundidad óptica de una losa delgada de altura $H(r)$ atravesada con ángulo de incidencia $\cos\theta = |\hat{d} \cdot \hat{n}|$ sigue la ley de Beer-Lambert:
\begin{equation}
\tau = \frac{H(r) \cdot \rho(r, \phi)}{\max(|\hat{d} \cdot \hat{n}|, \epsilon)}
\end{equation}
\begin{equation}
\alpha = 1 - e^{-\tau}
\end{equation}

\subsection{Efecto Doppler Relativista y Beaming}
El material del disco orbita con velocidad kepleriana $\vec{v}(r) = \Omega(r) r \hat{\phi}$, donde $\Omega(r) = \sqrt{\frac{GM}{r^3}}$.
El factor de velocidad relativista $\beta = \|\vec{v}\|/c$ determina el factor de Lorentz:
\begin{equation}
\gamma = \frac{1}{\sqrt{1 - \beta^2}}
\end{equation}
El corrimiento Doppler observado para una dirección de rayo $\hat{k} = -\hat{d}_{view}$ es:
\begin{equation}
\delta = \frac{1}{\gamma \left(1 - \vec{\beta} \cdot \hat{k}\right)}
\end{equation}

La intensidad observada sufre una amplificación por \textit{Doppler beaming} proporcional a la cuarta potencia de la frecuencia ($\delta^3$ a $\delta^4$):
\begin{equation}
I_{obs} = I_0 \cdot \delta^3
\end{equation}
Asimismo, el corrimiento al rojo gravitacional de la métrica atenúa la energía del fotón conforme escapa del pozo gravitatorio:
\begin{equation}
g = \sqrt{1 - \frac{r_s}{r}}
\end{equation}

\section{Turbulencia Multifractal y Cizalladura Kepleriana}

\subsection{Coordenadas en Espiral Logarítmica}
Dado que el fluido orbita diferencialmente según $\frac{d\Omega}{dr} < 0$, cualquier perturbación azimutal se estira en brazos espirales logarítmicos. La coordenada angular conforme deformada es:
\begin{equation}
\xi(r, \theta, t) = \theta + \Omega(r) t + k_s \ln\left(\frac{r}{r_{in}}\right)
\end{equation}
donde $k_s \approx 2.8 - 3.2$ es el factor de enrollamiento espiral.

\subsection{Micro-Estrías Filamentarias (MHD)}
Para reproducir las miles de corrientes de gas caliente visibles en el filme, se combinan armónicos de ruido procedural elípticos altamente alargados:
\begin{equation}
\mathcal{F}(\xi, r) = \sum_{j=1}^3 w_j \cdot \text{noise}\left(2^j \xi,\, f_j \frac{r - r_{in}}{r_{out} - r_{in}}\right)
\end{equation}
con frecuencias radiales $f_1 = 45$, $f_2 = 110$, $f_3 = 260$.

\section{Pipeline de Post-Procesamiento y Óptica Fílmica}

\subsection{Arquitectura Multi-Pase HDR}
El renderizador opera en espacio lineal de punto flotante de 16 bits por canal (\texttt{GL\_RGBA16F}):
\begin{enumerate}
    \item \textbf{Pase 1: Raymarching:} Generación de radiancia $L_{HDR}(\vec{x}) \in [0, \infty)$.
    \item \textbf{Pase 2: Extracción de Brillos (\textit{Bright Pass}):} Filtro con rodilla suave (\textit{soft-knee}):
    \begin{equation}
    f_{knee}(L) = \frac{\text{clamp}(L - L_{th} + k,\, 0,\, 2k)^2}{4k}
    \end{equation}
    \item \textbf{Pase 3: Bloom y Flare Anamórfico:} Desenfoque gaussiano separable $9$-tap colapsado a $5$-tap mediante muestreo bilineal de hardware con núcleos horizontales ensanchados:
    \begin{equation}
    \sigma_x \gg \sigma_y \implies \text{Anamorphic Horizontal Flare}
    \end{equation}
    \item \textbf{Pase 4: Composición y Mapeo Tonal ACES:}
    \begin{equation}
    f_{ACES}(x) = \frac{x (2.51 x + 0.03)}{x (2.43 x + 0.59) + 0.14}
    \end{equation}
\end{enumerate}

\section{Resultados y Rendimiento}

\begin{table}[htbp]
\centering
\caption{Métricas de Rendimiento en GPU (NVIDIA RTX @ 1080p)}
\label{tab:perf}
\begin{tabular}{lcc}
\toprule
\textbf{Fase del Pipeline} & \textbf{Tiempo ($\mu$s)} & \textbf{Presupuesto (\%)} \\
\midrule
Raymarching Escena (280 pasos) & 9400 & 82.5\% \\
Extracción de Brillos HDR & 180 & 1.6\% \\
Bloom Gaussiano (4 pases) & 520 & 4.6\% \\
Flare Anamórfico (5 pases) & 310 & 2.7\% \\
Composición ACES + Grano & 380 & 3.3\% \\
Overhead UI (Dear ImGui) & 600 & 5.3\% \\
\midrule
\textbf{Total por Frame} & \textbf{11.39 ms} & \textbf{87.8 FPS} \\
\bottomrule
\end{tabular}
\end{table}

\section{Conclusiones}
Se ha presentado un sistema riguroso e interactivo que combina los principios de la Relatividad General con técnicas avanzadas de computación gráfica en tiempo real. Mediante la correcta formulación de la transferencia radiativa, la inclusión de cizalla kepleriana multifractal y la preservación del pozo negro del horizonte de sucesos, se logra una reproducción fidedigna del agujero negro *Gargantua*.

\begin{thebibliography}{99}
\bibitem{thorne2015}
O. James, E. von Tunzelmann, P. Franklin, K. S. Thorne, \textit{"Gravitational lensing by spinning black holes in astrophysics, and in the movie Interstellar"}, Classical and Quantum Gravity, vol. 32, no. 6, 2015.

\bibitem{novikov1973}
I. D. Novikov, K. S. Thorne, \textit{"Astrophysics of Black Holes"}, in Black Holes (Les Astres Occlus), Gordon and Breach, New York, pp. 343--450, 1973.

\bibitem{shakura1973}
N. I. Shakura, R. A. Sunyaev, \textit{"Black holes in binary systems. Observational appearance"}, Astronomy and Astrophysics, vol. 24, pp. 337--355, 1973.

\bibitem{narkowicz2015}
K. Narkowicz, \textit{"ACES Filmic Tone Mapping Curve"}, Technical Blog, 2015.

\bibitem{bardeen1972}
J. M. Bardeen, W. H. Press, S. A. Teukolsky, \textit{"Rotating Black Holes: Locally Nonrotating Frames, Energy Extraction, and Scalar Green Functions"}, The Astrophysical Journal, vol. 178, pp. 347--370, 1972.
\end{thebibliography}

\end{document}
